\documentclass[10pt,a4paper]{amsart}
\usepackage[margin=19mm]{geometry}
\usepackage{amsmath,amssymb,mathtools}
\usepackage[scr=rsfs,cal=euler]{mathalfa}
\usepackage{xfrac}
\usepackage[unicode,hidelinks,bookmarks=false]{hyperref}
\newcommand{\ie}{i.e.\ }
\newcommand{\cf}{cf.\ }
\newcommand{\resp}{respectively\ }
\newcommand{\Subsection}{\subsection}
\providecommand{\nobreakdash}{\nobreak\hbox{-}}
\setlength{\emergencystretch}{2em}
\multlinegap0pt
\usepackage{bm}
\usepackage{bbm}
\usepackage{dsfont}
\usepackage{xspace}
\usepackage{cancel}
\usepackage{mathtools}
\usepackage{amsthm}
\makeatletter
\@ifundefined{lemma}{\newtheorem{lemma}{Lemma}}{}
\makeatother
\title[Resource-dependent process times in hybrid flexible flowshop]{Resource-dependent process times in hybrid flexible flowshops}
\author{Ioannis Avgerinos, Ioannis Mourtos, Dimitrios Papathanasiou, Georgios Zois}
\date{Source: arXiv:2510.18093v1 (CC BY 4.0)}
\begin{document}
\maketitle

\noindent\emph{Source and licence.} Resource-dependent process times in hybrid flexible flowshops. Version arXiv:2510.18093v1 (CC BY 4.0), distributed under the Creative Commons Attribution 4.0 International licence, as recorded in the arXiv licence metadata for this exact version and shown on the abstract page https://arxiv.org/abs/2510.18093v1. Full rights notice: licenses/arxiv-2510.18093-CC-BY-4.0.txt.\par\smallskip

\noindent\textit{Editorial scope.} Counted unit: the Lemma whose source label is \texttt{lem:proof} in the article (source0.tex lines 376--378, source label \texttt{lem:proof}), delivered together with the article's own complete original proof of it (source0.tex lines 379--430, one \texttt{proof} environment of 52 lines). No mathematics is written, repaired or re-derived: statement and proof are the article's own text line for line. Required context retained uncounted: eq:lp1 (lines 366--374); eq:lp2 (lines 366--374); eq:lp3 (lines 366--374). Every definition, display and label of those lines is retained unchanged. Omitted from this package and not redistributed: 1--365, 375--375, 431--820. Nothing inside a retained excerpt is omitted or altered: every retained body is delivered exactly as the article's lines are written, so no omission receipt of any kind accompanies this package. Citations: the retained excerpts carry no citation command at all, so no bibliography entry of the article is transcribed and no reference is redistributed. Reference adaptation: none was needed and none was made; every reference inside the delivered unit resolves inside it, so no unresolved reference or citation is produced. Printed slips kept literal: no printed word, symbol, hypothesis or number of the counted statement or of its proof was altered. Layout and class: the article's own class is \texttt{article} with packages including amsmath, amsthm, amssymb, graphicx, multirow, longtable, wrapfig, float, xcolor, pdflscape, hyperref, setspace, mathtools, algorithm; the wrapper is the lane's pinned \texttt{amsart} editorial wrapper at 10pt on A4, which restates the theorem-like environments the retained text needs and adds the article's own macro definitions that the retained text needs (none), with extra wrapper packages (bm, bbm, dsfont, xspace, cancel, mathtools). The delivered numbering is the unit's own. Assets: no retained span contains a figure environment, a graphics-inclusion command, a PDF-page inclusion, a table or a TikZ picture; no image file is copied into this package, the package declares no asset dependency and no image file is redistributed with it. Licence: version arXiv:2510.18093v1 of this article is distributed under the Creative Commons Attribution 4.0 International licence, as recorded in the arXiv licence metadata for this exact version and shown on the abstract page https://arxiv.org/abs/2510.18093v1. A full rights notice is delivered at licenses/arxiv-2510.18093-CC-BY-4.0.txt. Characters: every retained excerpt is pure ASCII, so the delivered engine, XeLaTeX with its default Latin Modern text font, reports no missing character.
        \begin{flalign}
            \text{(LP): }& &&\notag &&\\
            \text{min }&C_{max} &&\notag &&\\
            &\sum_{w\in W^{*}}x_{j^{\prime}w}\geq 1 &&\forall j^{\prime}\in J^{*} \label{eq:lp1}&&\\
            &C_{max}\geq \sum_{j^{\prime}\in J^{*}}p_{j^{\prime}1}\cdot x_{j^{\prime}w} &&\forall w\in W^{*} \label{eq:lp2} &&\\
            &C_{max}\geq p_{js^{i}_{j}w^{+}_{s}} &&\forall j\in J, i\in \{1, ..., |S_{j}|\} \label{eq:lp3} &&\\
            &x_{j^{\prime}w}\geq 0 &&\forall j\in J^{*}, w\in W^{*} \notag &&\\
            &C_{max}\geq 0 &&\notag &&
        \end{flalign}

        \begin{lemma}
            The optimal solution of (LP) provides a lower bound on the makespan of any feasible solution to the malleable job scheduling problem. \label{lem:proof}
        \end{lemma}

        \begin{proof}
            Constraints (\ref{eq:lp1}) ensure that each operation is fully assigned to workers, i.e., the total fraction of an operation assigned across all workers equals one.

            Property P.1 guarantees that Constraint (\ref{eq:lp3}) holds: when an operation $(j, s^{i}_{j})$ is split among multiple workers, its processing time decreases as the number of assigned workers increases, reaching its minimum when the operation is split among the maximum allowed number $w^{+}_{s}$. Hence, the makespan of any feasible schedule must be at least as large as the processing time corresponding to the maximum possible split of any operation.

            By property P.2, the product $|Q| \cdot f_{j'}(Q)$ is non-decreasing with $|Q|$, for any subset of workers $Q$, where $|Q|\geq 1$ denotes the number of workers assigned to operation $j'$. This implies that:
            \begin{equation}
                p_{j'|Q|} \geq \frac{p_{j'1}}{|Q|} \label{eq-work}
            \end{equation}
            
            \medskip
            
            Let us now consider a feasible schedule $\texttt{S}$ of the malleable scheduling problem, where each operation $j^{\prime}$ is assigned to a subset of workers $Q_{j^{\prime}}$ and let $\texttt{S}_w$ be the subset of operations assigned to each worker $w$, and $C^{\texttt{S}}_{max}$ be the makespan of $\texttt{S}$. It must hold that the total load of each worker is bounded above by $C^{\texttt{S}}_{max}$:
             \begin{flalign}
                &C^{\texttt{S}}_{max} \geq \sum_{j' \in \texttt{S}_w}p_{j'|Q_{j^{\prime}}|} && \forall w\in W^*  && && && && &&\label{eq-work1}
            \end{flalign}
            A feasible solution to (LP), denoted by $\bar{\texttt{S}}$, can be defined by setting, for each operation $j^{\prime}\in J^{*}$ and each worker $w\in Q_{j^{\prime}}$ in $\texttt{S}$:
            \[
            \bar{x}_{j'w} =
            \begin{cases}
            \frac{1}{|Q_{j^{\prime}}|}, & \text{if } x_{j^{\prime}w} = 1, \\
            0, & \text{otherwise}.
            \end{cases}
            \]
            
          By substituting solution $\bar{x}_{j'w}$ to Constraint (\ref{eq:lp2}), the total load of each worker $w$ in $\bar{\texttt{S}}$ becomes equal to $\sum_{j' \in \texttt{S}_{w}} \frac{p_{j'1}}{|Q_{j^{\prime}}|}$, and less or equal to the makespan, $C_{max}$, of $\bar{\texttt{S}}$. Note that, for a worker $w^{\prime}$, where the maximum load is achieved, $C_{max} = \sum_{j^{\prime}\in \bar{\texttt{S}}_{w^{\prime}}} \frac{p_{j'1}}{|Q_{j^{\prime}}|}$, where $\bar{\texttt{S}}_{w^\prime}$ is the subset of operations assigned to worker $w^{\prime}$ in $\bar{S}$. Applying (\ref{eq-work}), combined with (\ref{eq-work1}), we yield that: 
           \begin{flalign}
            &C^{\texttt{S}}_{max} \geq \sum_{j'\in \texttt{S}_{w^{\prime}}} p_{j'|Q_{j^{\prime}}|} \geq C_{max}, && \notag && && && && &&
        \end{flalign} 

        %    completion time $C_{w}$ of worker $w$ is $S$ satisfies:
         %   \begin{equation}
          %      C_{w} \geq \sum_{j' \in J^{*}} p_{j'1} \cdot x_{j'w} 
           %     \qquad \forall w \in W^{*} \notag
            %\end{equation}
            
       %     Since the makespan $C_{\max}$ is the maximum completion time among all workers, it follows that:
        %    \begin{equation}
         %       C_{\max} \ge \sum_{j' \in J^{*}} p_{j'1} \cdot x_{j'w}
          %      \qquad \forall w \in W^{*} \notag
           % \end{equation}
      %      This corresponds exactly to constraint (\ref{eq:lp2}).


   %   As shown above, for any operation $j^{\prime}$, the processing time is $p_{j'|Q_{j^{\prime}}|} \ge \frac{p_{j'1}}{|Q_{j^{\prime}}|}$. 
      
       Let $C^*_{max}$ optimal objective value of the (LP). Since $C_{max}$ is the value of a feasible solution to (LP), it should hold that:
          \begin{flalign}
            &C^{\texttt{S}}_{max}\geq C^{*}_{max}, && \notag && (\text{LB}_{8})&& && && &&
        \end{flalign}  
        for any feasible solution $\texttt{S}$ to the malleable job scheduling problem.
        \end{proof}

\end{document}
